\chapter{量化与低精度计算}

量化的目标，是用更少的位表示权重、激活或 KV Cache，从而减少内存占用、降低带宽压力，并使用更高吞吐的低精度指令。它不是简单地把 float 强制转换成 int8。真正的量化包含比例因子、零点、分组、校准、误差控制、低精度乘加和反量化融合。若这些环节处理不好，模型质量和性能都可能变差。

\section{仿射量化模型}

常见 int8 仿射量化把实数 $x$ 近似为 $x \approx s(q-z)$，其中 $q$ 是整数，$s$ 是 scale，$z$ 是 zero point。对称量化常令 $z=0$，实现简单；非对称量化能更好覆盖偏移分布，但计算更复杂。权重量化通常可以离线完成，激活量化可能需要运行时统计或校准。

\begin{lstlisting}[language=C++,caption={教学版对称 int8 量化和反量化}]
#include <algorithm>
#include <cmath>
#include <cstdint>
#include <span>

void quantize_symmetric_i8(std::span<const float> input,
                           std::span<std::int8_t> output,
                           float scale) {
    const float inv_scale = 1.0F / scale;
    for (std::int64_t i = 0; i < static_cast<std::int64_t>(input.size()); ++i) {
        const float scaled = std::round(input[static_cast<std::size_t>(i)] *
                                        inv_scale);
        const float clamped = std::clamp(scaled, -127.0F, 127.0F);
        output[static_cast<std::size_t>(i)] =
            static_cast<std::int8_t>(clamped);
    }
}

void dequantize_symmetric_i8(std::span<const std::int8_t> input,
                             std::span<float> output,
                             float scale) {
    for (std::int64_t i = 0; i < static_cast<std::int64_t>(input.size()); ++i) {
        output[static_cast<std::size_t>(i)] =
            static_cast<float>(input[static_cast<std::size_t>(i)]) * scale;
    }
}
\end{lstlisting}

scale 的选择决定误差。如果 scale 太大，小值会被压到同一个整数，精度差；如果 scale 太小，大值会饱和。常见校准方法会根据最大绝对值、百分位数或误差最小化选择范围。权重中若存在离群值，逐张量量化可能损失明显，分组量化或逐通道量化会更稳。

\section{低精度矩阵乘}

int8 矩阵乘通常把 int8 输入和权重相乘，累加到 int32，再乘 scale 得到 float 或低精度输出。这样做的原因是，多个 int8 乘积累加后范围会超过 int8。硬件可能提供点积指令，一条指令完成多个 int8 乘法并累加到 int32。性能提升来自两个方面：数据更小，带宽压力降低；低精度点积指令吞吐更高。

但量化 GEMM 的数据流更复杂。它需要读取 scale，处理 zero point，可能需要 per-channel 或 per-group 反量化，还要决定反量化发生在内核内部还是写回后。若每个输出元素都频繁读取 scale 或做复杂校正，收益会下降。因此高性能量化内核会把反量化、bias、激活尽量融合到输出阶段。

\begin{keyidea}
量化不是单纯压缩模型文件，而是改变整个性能模型。权重变小会减少带宽，低精度指令会提高计算峰值，但 scale、zero point、反量化和误差控制会引入新的成本。
\end{keyidea}

\section{bf16、fp16 与累加精度}

除了 int8，CPU 上还常见 bf16 和 fp16。bf16 保留 float32 的指数位范围，尾数更短，适合深度学习中范围大但精度要求相对可控的场景；fp16 尾数更长但指数范围更小。许多硬件支持 bf16/fp16 乘法并累加到 float32，以平衡吞吐和精度。

选择低精度格式时，要区分存储格式、计算格式和累加格式。权重可以用 int8 存储，计算时转成 int32 累加；激活可以用 bf16 存储，累加用 float32；KV Cache 可以量化存储，读取时按块反量化。不同组合对应不同误差和性能。

\section{量化粒度}

逐张量量化只用一个 scale，元数据少，计算简单，但对离群值敏感。逐通道量化为每个输出通道使用 scale，常用于权重。分组量化把一段连续权重共享一个 scale，是大语言模型权重量化常见方案。粒度越细，误差越小，但 scale 元数据越多，内核处理越复杂。

KV Cache 量化还要考虑动态增长和在线写入。每个 token、每个 head、每个 block 如何选择 scale，会影响读取效率和误差。长上下文 decode 中，KV Cache 占用巨大，量化它能显著降低内存流量，但必须验证对困惑度、生成质量和长上下文稳定性的影响。

\begin{exercisebox}
取一组 float 权重，分别用逐张量 scale 和每 32 个元素一组的 scale 做 int8 量化。计算反量化后的最大绝对误差和平均绝对误差。观察离群值对逐张量量化的影响。
\end{exercisebox}
